\section{DEP-001--003: Migrate from js-yaml to yaml v2}
\label{sec:dep-001-003}

\subsection*{Metadata}
\begin{tabular}{ll}
\textbf{Issue ID} & DEP-001, DEP-002, DEP-003 \\
\textbf{Severity} & Medium \\
\textbf{Category} & Dependencies \\
\textbf{Branch} & \branchref{fix/DEP-001-003-migrate-to-yaml-v2} \\
\textbf{Changes} & \branchdiff{fix/DEP-001-003-migrate-to-yaml-v2} \\
\end{tabular}

\subsection*{Executive Summary}
The project depended on two YAML parsing libraries simultaneously: \texttt{js-yaml} v4.1.0 and \texttt{yaml} v2.8.2. The \texttt{js-yaml} library is outdated with known security concerns, while \texttt{yaml} v2 is actively maintained and supports YAML 1.2. Maintaining two libraries doubles the dependency surface area, increases bundle size by approximately 750KB, and creates confusion about which library to use for new code.

The migration replaces all 53 instances of \texttt{js-yaml} \texttt{yaml.load()} and \texttt{yaml.dump()} calls with the \texttt{yaml} v2 equivalents (\texttt{YAML.parse()} and \texttt{YAML.stringify()}) across 19 source files. The \texttt{js-yaml} and \texttt{@types/js-yaml} dependencies were removed from \texttt{package.json}, reducing the install size by approximately 350KB.

\subsection*{Root Cause Analysis}

The source files imported both libraries side-by-side:

\begin{lstlisting}[language=TypeScript]
// src/plugin-cli.ts — Before fix: dual imports
import yaml from "js-yaml";           // Line 4
import { parseDocument } from "yaml"; // Line 7
\end{lstlisting}

\texttt{js-yaml} was used for quick \texttt{yaml.load()} operations on agent manifests and plugin lists across 19 files. \texttt{yaml} v2 was imported specifically for \texttt{parseDocument()} in \texttt{plugin-cli.ts} because it preserves comments and formatting when editing \texttt{agent.yaml}---something \texttt{js-yaml} does not support.

The key API differences:

\begin{tabular}{lll}
\textbf{Operation} & \textbf{js-yaml} & \textbf{yaml v2} \\
\hline
Parse YAML string & \texttt{yaml.load(str)} & \texttt{YAML.parse(str)} \\
Serialize to YAML & \texttt{yaml.dump(obj)} & \texttt{YAML.stringify(obj)} \\
Comment-preserving parse & Not supported & \texttt{parseDocument(str)} \\
\end{tabular}

\subsection*{Fix Implementation}

All 19 source files were updated to use \texttt{yaml} v2 exclusively. The changes followed a consistent pattern across every file:

\begin{lstlisting}[language=TypeScript]
// Before: js-yaml import
import yaml from "js-yaml";

// After: yaml v2 import
import yaml from "yaml";
\end{lstlisting}

API calls were updated from \texttt{js-yaml} to \texttt{yaml} v2 syntax:

\begin{lstlisting}[language=TypeScript]
// Before (js-yaml)
const manifest = yaml.load(raw) as AgentManifest;
const yamlStr = yaml.dump(frontmatter, { lineWidth: -1, noRefs: true });

// After (yaml v2)
const manifest = yaml.parse(raw) as AgentManifest;
const yamlStr = yaml.stringify(frontmatter, { lineWidth: -1 });
\end{lstlisting}

Dependencies were updated in \texttt{package.json}:

\begin{lstlisting}[language=diff]
-    "js-yaml": "^4.1.0",
-    "@types/js-yaml": "^4.0.9",
\end{lstlisting}

\subsection*{Affected Files}
All 19 files that were migrated:
\begin{itemize}
    \item \srcfile{package.json} --- removed \texttt{js-yaml} and \texttt{@types/js-yaml}
    \item \srcfile{src/agents.ts} --- frontmatter parsing
    \item \srcfile{src/compliance.ts} --- regulatory config loading
    \item \srcfile{src/config.ts} --- environment config loading
    \item \srcfile{src/hooks.ts} --- hooks configuration
    \item \srcfile{src/knowledge.ts} --- knowledge index loading
    \item \srcfile{src/learning/reinforcement.ts} --- frontmatter parsing with \texttt{stringify}
    \item \srcfile{src/loader.ts} --- agent manifest loading
    \item \srcfile{src/plugin-cli.ts} --- plugin list and agent.yaml editing
    \item \srcfile{src/plugins.ts} --- plugin manifest loading
    \item \srcfile{src/schedules.ts} --- schedule definitions with \texttt{stringify}
    \item \srcfile{src/skills.ts} --- skill metadata parsing
    \item \srcfile{src/tool-loader.ts} --- declarative tool definitions
    \item \srcfile{src/tools/memory.ts} --- memory config loading
    \item \srcfile{src/tools/sandbox-memory.ts} --- sandbox memory config
    \item \srcfile{src/tools/skill-learner.ts} --- skill frontmatter with \texttt{stringify}
    \item \srcfile{src/tools/task-tracker.ts} --- task skill searching
    \item \srcfile{src/workflows.ts} --- workflow definitions with \texttt{stringify}
    \item \srcfile{src/\_\_tests\_\_/dep-001-003.test.ts} --- migration verification tests
\end{itemize}

\subsection*{Verification}
\begin{lstlisting}[language=TypeScript]
// dep-001-003.test.ts
import { parse, stringify, parseDocument } from "yaml";

// Test 1: Basic YAML parse works
const basic = parse("name: test\nversion: 1.0\n");
assert.strictEqual(basic.name, "test");
assert.strictEqual(basic.version, 1.0);

// Test 2: Plugin manifest parsing
const manifestYaml = `
id: my-plugin
name: My Plugin
version: 0.1.0
description: Test
provides:
  tools: true
  skills: true
`.trim();
const manifest = parse(manifestYaml);
assert.strictEqual(manifest.id, "my-plugin");
assert.strictEqual(manifest.provides.tools, true);

// Test 3: Round-trip stringify preserves structure
const obj = { a: 1, b: { c: [1, 2, 3] } };
const yml = stringify(obj);
const parsed = parse(yml);
assert.deepStrictEqual(parsed, obj);

// Test 4: agent.yaml editing preserves comments
const input = `# Agent configuration
name: my-agent
version: 1.0.0
# Model settings
model:
  preferred: "anthropic:claude-4"
`.trim();
const doc = parseDocument(input);
doc.set("version", "2.0.0");
const output = doc.toString();
assert.ok(output.includes("# Agent configuration"));
assert.ok(output.includes("# Model settings"));
assert.ok(output.includes('version: "2.0.0"'));

// Test 5: No js-yaml dependency remains
const pkg = JSON.parse(readFileSync("./package.json", "utf-8"));
const deps = { ...pkg.dependencies, ...pkg.devDependencies };
ok(!deps["js-yaml"], "js-yaml removed from dependencies");
ok(!deps["@types/js-yaml"], "@types/js-yaml removed from devDependencies");
\end{lstlisting}
